\documentclass[10pt,a4paper]{amsart}
\usepackage{amsmath,amssymb,amsthm,graphicx}
\usepackage{verbatim}
\newenvironment{reference}{\comment}{\endcomment}
%\newenvironment{reference}{}{}
\newenvironment{slogan}{\comment}{\endcomment}
\newenvironment{history}{\comment}{\endcomment}
% For commutative diagrams we use Xy-pic
\usepackage[all]{xy}

% We use 2cell for 2-commutative diagrams.
\xyoption{2cell}
% Theorem environments.
%
\theoremstyle{plain}
\newtheorem{theorem}[subsection]{Theorem}
\newtheorem{proposition}[subsection]{Proposition}
\newtheorem{lemma}[subsection]{Lemma}

\theoremstyle{definition}
\newtheorem{definition}[subsection]{Definition}
\newtheorem{example}[subsection]{Example}
\newtheorem{exercise}[subsection]{Exercise}
\newtheorem{situation}[subsection]{Situation}

\theoremstyle{remark}
\newtheorem{remark}[subsection]{Remark}
\newtheorem{remarks}[subsection]{Remarks}

\numberwithin{equation}{subsection}

% Macros
%
\def\lim{\mathop{\mathrm{lim}}\nolimits}
\def\colim{\mathop{\mathrm{colim}}\nolimits}
\def\Spec{\mathop{\mathrm{Spec}}}
\def\Hom{\mathop{\mathrm{Hom}}\nolimits}
\def\Ext{\mathop{\mathrm{Ext}}\nolimits}
\def\SheafHom{\mathop{\mathcal{H}\!\mathit{om}}\nolimits}
\def\SheafExt{\mathop{\mathcal{E}\!\mathit{xt}}\nolimits}
\def\Sch{\mathit{Sch}}
\def\Mor{\mathop{\mathrm{Mor}}\nolimits}
\def\Ob{\mathop{\mathrm{Ob}}\nolimits}
\def\Sh{\mathop{\mathit{Sh}}\nolimits}
\def\NL{\mathop{N\!L}\nolimits}
\def\CH{\mathop{\mathrm{CH}}\nolimits}
\def\proetale{{pro\text{-}\acute{e}tale}}
\def\etale{{\acute{e}tale}}
\def\QCoh{\mathit{QCoh}}
\def\Ker{\mathop{\mathrm{Ker}}}
\def\Im{\mathop{\mathrm{Im}}}
\def\Coker{\mathop{\mathrm{Coker}}}
\def\Coim{\mathop{\mathrm{Coim}}}

% Boxtimes
%
\DeclareMathSymbol{\boxtimes}{\mathbin}{AMSa}{"02}

%
% Macros for moduli stacks/spaces
%
\def\QCohstack{\mathcal{QC}\!\mathit{oh}}
\def\Cohstack{\mathcal{C}\!\mathit{oh}}
\def\Spacesstack{\mathcal{S}\!\mathit{paces}}
\def\Quotfunctor{\mathrm{Quot}}
\def\Hilbfunctor{\mathrm{Hilb}}
\def\Curvesstack{\mathcal{C}\!\mathit{urves}}
\def\Polarizedstack{\mathcal{P}\!\mathit{olarized}}
\def\Complexesstack{\mathcal{C}\!\mathit{omplexes}}
% \Pic is the operator that assigns to X its picard group, usage \Pic(X)
% \Picardstack_{X/B} denotes the Picard stack of X over B
% \Picardfunctor_{X/B} denotes the Picard functor of X over B
\def\Pic{\mathop{\mathrm{Pic}}\nolimits}
\def\Picardstack{\mathcal{P}\!\mathit{ic}}
\def\Picardfunctor{\mathrm{Pic}}
\def\Deformationcategory{\mathcal{D}\!\mathit{ef}}
\usepackage{cleveref}
\crefname{originallemma}{original lemma}{original lemmas}
\Crefname{originallemma}{Original lemma}{Original lemmas}
\makeatletter
\newlabel{lemma-the-complex-for-L-star}{{0FUF}{source}{}{sourcepointer-1}{}}
\newlabel{lemma-the-complex-for-L-star@cref}{{[originallemma][0][]0FUF}{[1][0][]source}{}{}{}}
\makeatother
\crefname{originalremark}{original remark}{original remarks}
\Crefname{originalremark}{Original remark}{Original remarks}
\makeatletter
\newlabel{remark-cup-product-as-a-map}{{0FU7}{source}{}{sourcepointer-2}{}}
\newlabel{remark-cup-product-as-a-map@cref}{{[originalremark][0][]0FU7}{[1][0][]source}{}{}{}}
\makeatother
\crefname{originalsection}{original section}{original sections}
\Crefname{originalsection}{Original section}{Original sections}
\makeatletter
\newlabel{sdga-section-misc}{{0FU3}{source}{}{sourcepointer-3}{}}
\newlabel{sdga-section-misc@cref}{{[originalsection][0][]0FU3}{[1][0][]source}{}{}{}}
\makeatother
\crefname{originalsection}{original section}{original sections}
\Crefname{originalsection}{Original section}{Original sections}
\makeatletter
\newlabel{section-first-chern-class}{{0FLE}{source}{}{sourcepointer-4}{}}
\newlabel{section-first-chern-class@cref}{{[originalsection][0][]0FLE}{[1][0][]source}{}{}{}}
\makeatother
\title{The boundary map of the punctured-line-bundle sequence is the map attached to the first de Rham Chern class}
\author{The Stacks Project Authors}
\date{stacks-project:derham.tex:local-snapshot-20261012; GNU FDL 1.2 or later}
\begin{document}
\maketitle
\noindent\textit{Editorial note.} All delivered mathematical text is the author's own, copied from the pinned Stacks Project chapter file derham.tex; no statement or proof step was rewritten, reordered or removed.\par\medskip
\noindent\textit{Editorial note.} Preamble provenance: the pinned chapter file opens with \textbackslash{}input\{preamble\} and carries no preamble text of its own, so the project's single Stacks-wide preamble staged at sources/stacks-derham/preamble.tex is excerpted verbatim as three spans, lines 8-12 (verbatim package and the project's own reference/slogan/history display environments), lines 14-18 (Xy-pic setup) and lines 154-217 (theorem declarations and mathematical macros). The wrapper class and the amsmath/amssymb/amsthm/graphicx base packages are supplied by the assembler. No macro was invented or redefined.\par\medskip
\noindent\textit{Editorial note.} Same-chapter pointers are located at their \textbackslash{}label lines inside the pinned chapter file; cross-chapter pointers are located in the project's own tag list staged in the snapshot (sources/stacks-chow/tags.txt), at the verbatim lines carrying their canonical tags.\par\medskip
\noindent\textit{Allowed standard prerequisites (editorial).}\begin{itemize}
\item The short exact sequence of the punctured line bundle and the residue map. Same chapter derham.tex, label lemma-the-complex-for-L-star (invoked at proof lines 1426, 1438, 1452 and 1491). It provides the shifted sequence \$0 \textbackslash{}to \textbackslash{}Omega\textasciicircum{}\textbackslash{}bullet\_\{X/S\}[1] \textbackslash{}to \textbackslash{}Omega\textasciicircum{}\textbackslash{}bullet\_\{L\textasciicircum{}\textbackslash{}star/S,0\}[1] \textbackslash{}to \textbackslash{}Omega\textasciicircum{}\textbackslash{}bullet\_\{X/S\} \textbackslash{}to 0\$ whose boundary is \$\textbackslash{}delta\$, and the local description of the residue map.
\item The map attached to a de Rham cohomology class, and its uniqueness. Same chapter derham.tex, remark remark-cup-product-as-a-map (invoked at proof lines 1427 and 1459). It defines the map of the statement for a class \$\textbackslash{}xi\$ and states its uniqueness, which reduces the identification to checking the value on \$1\$.
\item Right differential graded modules over the de Rham complex and the derived category computation. Differential Graded Sheaves, Section \textbackslash{}ref\{sdga-section-misc\} (invoked at proof line 1454). It shows that the boundary map, being a map of right differential graded \$\textbackslash{}Omega\textasciicircum{}\textbackslash{}bullet\_\{X/S\}\$-modules, comes from a map in the derived category of such modules.
\item Local trivializations and the Cech cocycle of an invertible sheaf. Same chapter derham.tex, Section \textbackslash{}ref\{section-first-chern-class\} (invoked at proof line 1495). They present the class of \$\textbackslash{}mathcal\{L\}\$ by the cocycle \$\textbackslash{}\{f\_\{i\_0i\_1\}\textbackslash{}\}\$, which is what the final computation of the boundary uses.
\end{itemize}\medskip
\ifdefined\hypertarget\hypertarget{sourcepointer-1}{}\fi
\noindent\textit{Source reference.} Original lemma 0FUF: sources/stacks-derham/pinned-original-source.tex, source lines 1348--1348. Navigation pointer for \textbackslash{}ref\{lemma-the-complex-for-L-star\} used inside the counted span (line 1426); the argument that consumes it is retained in the span.\par\medskip
\ifdefined\hypertarget\hypertarget{sourcepointer-2}{}\fi
\noindent\textit{Source reference.} Original remark 0FU7: sources/stacks-derham/pinned-original-source.tex, source lines 575--575. Navigation pointer for \textbackslash{}ref\{remark-cup-product-as-a-map\} used inside the counted span (line 1427); the argument that consumes it is retained in the span.\par\medskip
\ifdefined\hypertarget\hypertarget{sourcepointer-3}{}\fi
\noindent\textit{Source reference.} Original section 0FU3: sources/stacks-chow/tags.txt, source lines 19763--19763. Navigation pointer for \textbackslash{}ref\{sdga-section-misc\} used inside the counted span (line 1454); the argument that consumes it is retained in the span.\par\medskip
\ifdefined\hypertarget\hypertarget{sourcepointer-4}{}\fi
\noindent\textit{Source reference.} Original section 0FLE: sources/stacks-derham/pinned-original-source.tex, source lines 1035--1035. Navigation pointer for \textbackslash{}ref\{section-first-chern-class\} used inside the counted span (line 1495); the argument that consumes it is retained in the span.\par\medskip
\begin{lemma}
\label{lemma-the-complex-for-L-star-gives-chern-class}
The ``boundary'' map
$\delta : \Omega^\bullet_{X/S} \to \Omega^\bullet_{X/S}[2]$
in $D(X, f^{-1}\mathcal{O}_S)$ coming from
the short exact sequence in Lemma \ref{lemma-the-complex-for-L-star}
is the map of Remark \ref{remark-cup-product-as-a-map}
for $\xi = c_1^{dR}(\mathcal{L})$.
\end{lemma}

\begin{proof}
To be precise we consider the shift
$$
0 \to \Omega^\bullet_{X/S}[1] \to
\Omega^\bullet_{L^\star/S, 0}[1] \to
\Omega^\bullet_{X/S} \to 0
$$
of the short exact sequence of Lemma \ref{lemma-the-complex-for-L-star}.
As the degree zero part of a grading on
$(L^\star \to X)_*\Omega^\bullet_{L^\star/S}$
we see that $\Omega^\bullet_{L^\star/S, 0}$ is a differential
graded $\mathcal{O}_X$-algebra and that the map
$\Omega^\bullet_{X/S} \to \Omega^\bullet_{L^\star/S, 0}$
is a homomorphism of differential graded $\mathcal{O}_X$-algebras.
Hence we may view $\Omega^\bullet_{X/S}[1] \to
\Omega^\bullet_{L^\star/S, 0}[1]$ as a map of right differential graded
$\Omega^\bullet_{X/S}$-modules on $X$. The map
$\text{Res} : \Omega^\bullet_{L^\star/S, 0}[1] \to \Omega^\bullet_{X/S}$
is a map of right differential graded $\Omega^\bullet_{X/S}$-modules
since it is locally defined by the rule
$\text{Res}(\omega' + \text{d}\log(s) \wedge \omega) = \omega$, see
proof of Lemma \ref{lemma-the-complex-for-L-star}.
Thus by the discussion in
Differential Graded Sheaves, Section \ref{sdga-section-misc}
we see that $\delta$ comes from a map
$\delta' : \Omega^\bullet_{X/S} \to \Omega^\bullet_{X/S}[2]$
in the derived category $D(\Omega^\bullet_{X/S}, \text{d})$
of right differential graded modules over the de Rham complex.
The uniqueness averted in Remark \ref{remark-cup-product-as-a-map}
shows it suffices to prove that $\delta(1) = c_1^{dR}(\mathcal{L})$.

\medskip\noindent
We claim that there is a commutative diagram
$$
\xymatrix{
0 \ar[r] &
\mathcal{O}_X^* \ar[r] \ar[d]_{\text{d}\log} &
E \ar[r] \ar[d] &
\underline{\mathbf{Z}} \ar[d] \ar[r] &
0 \\
0 \ar[r] &
\Omega^\bullet_{X/S}[1] \ar[r] &
\Omega^\bullet_{L^\star/S, 0}[1] \ar[r] &
\Omega^\bullet_{X/S} \ar[r] &
0
}
$$
where the top row is a short exact sequence of abelian sheaves whose
boundary map sends $1$ to the class of $\mathcal{L}$ in
$H^1(X, \mathcal{O}_X^*)$. It suffices to prove the claim
by the compatibility of boundary maps with maps between short
exact sequences. We define $E$ as the sheafification of the rule
$$
U \longmapsto \{(s, n) \mid
n \in \mathbf{Z},\ s \in \mathcal{L}^{\otimes n}(U)\text{ generator}\}
$$
with group structure given by $(s, n) \cdot (t, m) = (s \otimes t, n + m)$.
The middle vertical map sends $(s, n)$ to $\text{d}\log(s)$. This produces
a map of short exact sequences
because the map $Res : \Omega^1_{L^\star/S, 0} \to \mathcal{O}_X$
constructed in the proof of Lemma \ref{lemma-the-complex-for-L-star} sends
$\text{d}\log(s)$ to $1$ if $s$ is a local generator of $\mathcal{L}$.
To calculate the boundary of $1$ in the top row, choose local trivializations
$s_i$ of $\mathcal{L}$ over opens $U_i$ as in
Section \ref{section-first-chern-class}. On the overlaps
$U_{i_0i_1} = U_{i_0} \cap U_{i_1}$
we have an invertible function $f_{i_0i_1}$ such that
$f_{i_0i_1} = s_{i_1}|_{U_{i_0i_1}} s_{i_0}|_{U_{i_0i_1}}^{-1}$
and the cohomology class of $\mathcal{L}$ is given by the {\v C}ech cocycle
$\{f_{i_0i_1}\}$. Then of course we have
$$
(f_{i_0i_1}, 0) = (s_{i_1}, 1)|_{U_{i_0i_1}} \cdot
(s_{i_0}, 1)|_{U_{i_0i_1}}^{-1}
$$
as sections of $E$ which finishes the proof.
\end{proof}

\end{document}
